\section{Xây dựng Cloud Server}

\subsection{Điểm khởi chạy và vòng đời}

Tệp \path{apps/cloud/main.py} tạo ứng dụng FastAPI từ hai container phụ thuộc:
\texttt{CloudContainer} (repository, service, handler, kết nối MQTT) và
\texttt{RealtimeContainer} (luồng WebSocket). Trong \texttt{lifespan} của ứng
dụng, hệ thống lần lượt: khởi động container; chạy session manager của MCP trong
một task riêng; nếu bật embedded worker thì đăng ký các consumer MQTT và chạy
NodeLivenessService. Khi tắt, các task được hủy theo thứ tự ngược lại để không
rò rỉ kết nối.

Router REST được gắn tiền tố \texttt{/api/v1} và phụ thuộc \texttt{ApiKeyAuth};
ứng dụng MCP được \texttt{mount} tại \texttt{/mcp} và bọc bởi \texttt{ApiKeyGuard}
--- một ASGI middleware kiểm tra header \texttt{x-api-key}, vì ứng dụng được mount
không đi qua dependency của FastAPI. \texttt{/health} kiểm tra đồng thời database,
Redis và broker.

\subsection{REST API}

Bảng~\ref{tab:cloud-api} liệt kê các endpoint chính; đặc tả đầy đủ ở
Phụ lục~\ref{app:api}.

\begin{table}[H]
\centering
\caption{REST API của Cloud}
\label{tab:cloud-api}
\begin{tabularx}{\textwidth}{L{2.2cm}L{6.4cm}Y}
\toprule
\textbf{Method} & \textbf{Endpoint} & \textbf{Chức năng}\\
\midrule
GET/POST & \path{/api/v1/clusters} & Liệt kê, tạo cluster\\
GET/POST & \path{/api/v1/nodes} & Liệt kê (lọc theo cluster), đăng ký node\\
GET & \path{/api/v1/agents/{node_id}} & Trạng thái Agent kèm runtime\\
GET & \path{/api/v1/runtimes} & Liệt kê runtime (lọc theo node)\\
POST & \path{/api/v1/runtimes/{id}/{action}} & start/stop/restart\\
POST/GET & \path{/api/v1/deployments} & Tạo (202), liệt kê\\
PUT/DELETE & \path{/api/v1/deployments/{id}} & Cập nhật, gỡ bỏ (202)\\
GET & \path{/api/v1/monitoring/nodes/{id}/telemetry} & Telemetry mới nhất\\
GET & \path{/api/v1/monitoring/nodes/{id}/events} & Event gần đây\\
WS & \path{/api/v1/ws?topics=...} & Luồng trạng thái/telemetry/event\\
\bottomrule
\end{tabularx}
\end{table}

Lỗi nghiệp vụ được biểu diễn bằng \texttt{AppException} với mã chuẩn hóa
(\texttt{NOT\_FOUND}, \texttt{CONFLICT}, \texttt{INVALID\_REQUEST},
\texttt{EXTERNAL\_SERVICE}...) và được \texttt{AppExceptionHandler} chuyển thành
mã HTTP tương ứng, giúp client và Assistant xử lý lỗi thống nhất.

\subsection{Tầng lưu trữ}

Repository PostgreSQL dùng \texttt{asyncpg} với truy vấn SQL viết tay và lớp
\texttt{RowMapper} chuyển bản ghi thành entity. Schema gồm hai migration: V001
tạo các bảng người dùng, API key (chỉ lưu hàm băm và tiền tố hiển thị), cluster,
node, khóa API của node, model, lớp của model, runtime, nhật ký cluster/node/runtime
và phiên chat; V002 thêm bảng \texttt{deployments} và trạng thái DEGRADED cho
runtime. Bảng \texttt{deployments} cố ý không có khóa ngoại tới \texttt{runtimes}
vì bản ghi runtime bị xóa khi gỡ bỏ còn deployment được giữ làm lịch sử. Trigger
\texttt{set\_updated\_at} tự cập nhật thời điểm sửa đổi. Mỗi repository có một bản
hiện thực trong bộ nhớ tương ứng dùng cho kiểm thử.

\subsection{Tiêu thụ MQTT}

\texttt{CloudContainer.subscribe\_consumers()} đăng ký các handler ở
Mục~\ref{sec:cloud-design} trên ba wildcard state, telemetry, event. Mỗi
handler kiểm tra schema bằng Pydantic trước khi ghi; thông điệp lỗi chỉ bị ghi
log cảnh báo, không làm dừng consumer. Chú ý kỹ thuật: thư viện paho-mqtt cần
\texttt{add\_reader} của vòng sự kiện, nên trên Windows ứng dụng chủ động dùng
\texttt{SelectorEventLoop} thay cho Proactor mặc định.

\subsection{Kho model}

Khi biến môi trường \texttt{CLOUD\_MODEL\_STORE} có giá trị
\texttt{huggingface}, Cloud khởi tạo repository kho model
(\texttt{HuggingFaceModelRepository}) với token và endpoint được cấu hình. Repository hỗ
trợ \texttt{resolve} (phân giải revision thành commit SHA và URL tải), \texttt{upload}
(đẩy tệp model, trả về commit mới) và \texttt{ping}. Lời gọi thư viện đồng bộ được
đưa vào thread (\texttt{asyncio.to\_thread}) để không chặn API. Deployment tham
chiếu model bằng \texttt{model\_uri} là URL đã ghim, Edge tải trực tiếp từ kho model
nên băng thông tải model không đi qua Cloud Server.
